\chapter{Discussion and Contributions}

\section{Comparison with Existing TCM AI Systems}

TCM-Sage occupies a distinct position in the TCM AI landscape. While existing products pursue an ``AI Doctor'' paradigm, TCM-Sage functions as an evidence-synthesis reference tool. This distinction is not merely semantic; it reflects fundamentally different design philosophies with practical consequences.

General-purpose LLMs such as ChatGPT and Qwen operate as black boxes. When queried about TCM topics, they generate responses from parametric memory without any mechanism for source verification. A practitioner receiving a formula recommendation from ChatGPT cannot determine whether that recommendation originates from \textit{Shanghan Lun}, from a modern textbook, or from the model's interpolation of training data. In clinical contexts where incorrect dosage or formula selection carries real risk, this unverifiability is the primary obstacle to adoption, not hallucination per se.

Specialized TCM AI products have attempted to address domain accuracy but not transparency. HuatuoGPT achieved state-of-the-art results on multiple medical benchmarks \cite{huatuogpt2023}. Qihuang 2.0 (Shuzhi Qihuang), a 32-billion-parameter multimodal model \cite{qihuang2024ecnu}, scored 88.1\% on the TCM licensing examination and incorporated tongue and pulse image analysis. iFlytek's TCM assistant was deployed across primary care institutions in Anhui Province \cite{iflytek2024cnr}. These systems share a common ``AI Doctor'' design philosophy, focusing on direct diagnosis and prescription. TCM-Sage takes a fundamentally different direction: evidence synthesis with traceable citations from classical texts.

TCM-Sage addresses these barriers through five design decisions:

\begin{enumerate}
    \item \textbf{Citation Traceability}: Every claim links to a specific classical text passage, visible through the Citation Panel. Practitioners can verify sources without copy-pasting text into external search tools.
    \item \textbf{Classical Text Specialization}: The 17-text corpus provides depth in foundational TCM theory that general LLMs cannot match, backed by clause-level retrieval for structured texts.
    \item \textbf{Full Local Deployment}: The entire pipeline, including embedding models and LLM inference, can run on consumer hardware via Ollama or LMStudio. No patient data leaves the institution's network.
    \item \textbf{Modular Knowledge Base}: Adding or removing knowledge sources requires only text ingestion, not model retraining. When practitioners suggested adding modern clinical texts during testing, the response was immediate: the architecture supports this without modification.
    \item \textbf{Ancient Measurement Accuracy}: An embedded conversion table provides historically verified unit conversions (Han dynasty \textit{liang} $\approx$ 13.8g, not the commonly misquoted 3g), preventing dosage calculation errors that other systems produce.
\end{enumerate}

The SymMap v2.0 knowledge graph deserves specific mention. Unlike the classical texts, SymMap is a modern pharmacological database integrating contemporary research. The crosswalk bridge connects these modern entities to classical terminology, creating a bidirectional link between historical medical knowledge and current pharmacological understanding. This hybrid of classical texts and modern knowledge graph is unique among TCM AI systems.

\section{Strengths Demonstrated}

The Arena evaluation provides quantitative evidence: with $p = 0.0011$ and Cohen's $d = 0.45$, the preference for TCM-Sage over the baseline is statistically significant at the 1\% level with a medium effect size. This result was observed across five days of testing with different tester groups.

Qualitative feedback reinforced specific strengths. A TCM student consultant summarized four key advantages observed during extended testing:

\begin{enumerate}
    \item \textbf{Focused, relevant responses:} The system provides adequately elaborated answers with valid arguments while avoiding irrelevant details. Unlike general LLMs that tend to produce verbose, unfocused responses mixing TCM with Western medicine or unrelated topics, TCM-Sage stays within the classical framework.
    \item \textbf{Credible source extraction and integration:} The system extracts and integrates credible sources directly from TCM classics, presenting them alongside the generated response through the Citation Panel.
    \item \textbf{TCM-perspective clinical guidance:} The system guides practitioners in thinking from TCM perspectives, including \textit{Materia Medica} recommendations and treatment approach demonstrations grounded in classical methodology (\textit{Bian Zheng Lun Zhi}).
    \item \textbf{Practitioner-informed tuning:} The system prompt and retrieval pipeline were iteratively refined based on feedback from TCM students, resulting in output that aligns with how practitioners expect clinical information to be structured.
\end{enumerate}

The Citation Panel received the strongest individual endorsement, with testers describing it as eliminating the need to manually verify AI-generated claims against source texts. In one case, the system recommended \textit{Jiegeng Tang} (Platycodon Decoction) for a sore throat that the consulting student had not considered. Upon review, he confirmed this as a valid and potentially superior approach from a classical perspective, demonstrating the system's ability to surface overlooked classical formulas.

The modular architecture received indirect validation when practitioners at Guangdong Provincial Hospital of TCM suggested expanding the corpus to include modern clinical masters' works (Zhang Xichun's \textit{Yi Xue Zhong Zhong Can Xi Lu}, Shi Xuemin's acupuncture texts, and Dong's Extraordinary Points system). The fact that such expansion is architecturally straightforward confirms that the design decision to use RAG over fine-tuning was sound.

\section{Limitations and Honest Assessment}

Several limitations emerged during testing and must be acknowledged.

\textbf{Knowledge Boundary Errors.} The system's accuracy is bounded by its corpus. When a practitioner asked about \textit{Mahuang Tang} dosage, the system stated that one Han dynasty \textit{liang} equals approximately 3 grams, a significant error (the correct value is 13.8--15.6g based on archaeological evidence). This occurred because the classical texts contain dosage ratios but not modern unit conversions. The plain LLM baseline, with access to web search, retrieved the correct conversion. This limitation was addressed by embedding a historically verified measurement conversion table directly into the system prompt.

\textbf{Corpus Coverage Gaps.} The 17-text corpus covers foundational theory, \textit{Shanghan}/\textit{Wenbing}, herbology, formulary, and acupuncture. It does not cover surgery, orthopedics, ophthalmology, specialized gynecology or pediatrics texts, or pulse diagnosis monographs. Practitioner feedback was direct: ``These texts are too traditional. Real clinical learning requires modern clinical masters' experience and knowledge systems.'' This feedback, while critical, validates the system's modular design as the appropriate architecture for iterative corpus expansion.

\textbf{Diagnostic Over-Assertion.} The system initially presented diagnostic interpretations with excessive confidence. When analyzing a pediatric case, it concluded \textit{Pi Xu} (spleen deficiency) without sufficient evidence in the case data, and classified an initial rash presentation as \textit{Re Yu Yu Ying} (heat trapped in the nutritive level) when the clinical presentation better supported \textit{Xie Qi Zai Biao} (pathogen at the exterior). Hedging principles were added to the system prompt to mitigate this tendency, but the underlying challenge of calibrating LLM confidence in medical reasoning remains partially unresolved.

\textbf{Evaluation Scale.} The 59-vote Arena sample, while statistically sufficient for significance testing, represents a limited evaluator pool. Most voters were TCM students rather than senior practitioners. The Arena format measures overall preference but does not decompose quality into specific dimensions such as citation accuracy, diagnostic reasoning, or readability.

\section{Ethical Considerations}

TCM-Sage is positioned as an evidence reference tool, not a diagnostic system. This distinction carries ethical weight. The system does not claim to diagnose, and its outputs carry an explicit disclaimer that all suggestions are for clinical reference only, with final diagnostic and prescriptive authority belonging to licensed practitioners.

Data privacy is addressed architecturally. Full local deployment eliminates the need to transmit patient data to external servers. During testing, practitioners at Guangdong Provincial Hospital of TCM noted this as a meaningful advantage, given institutional policies restricting patient data sharing with cloud AI services.

The use of AI tools in developing this project is disclosed in the Acknowledgments section, in compliance with university policy.

\begin{table}[H]
\centering
\caption{Feature comparison between TCM-Sage and existing TCM AI systems.}
\label{tab:comparison}
\begin{tabularx}{\textwidth}{lXXXX}
\toprule
\textbf{Feature} & \textbf{TCM-Sage} & \textbf{HuatuoGPT} & \textbf{Qihuang 2.0} & \textbf{iFlytek TCM} \\
\midrule
Primary Approach & Evidence synthesis & AI Doctor (diagnosis) & AI Doctor (multimodal) & AI Doctor (primary care) \\
Citation Traceability & Clause-level, click-through & None reported & None reported & None reported \\
Knowledge Base & 17 classical texts + SymMap KG & Fine-tuned on medical corpus & TCM KG + tongue/pulse images & Expert KG + clinical data \\
Local Deployment & Full local (LLM + embedding) & Cloud only & Cloud only & Cloud only \\
Knowledge Update & Add text files, re-ingest & Retrain model & Retrain model & Retrain model \\
Evaluation & Arena blind A/B ($p=0.0011$) & Medical benchmark SOTA & TCM licensing exam (88.1\%) & 9,800+ assisted diagnoses \\
\bottomrule
\end{tabularx}
\end{table}
